% ECONOMICS_PROBLEM: {"id": "C72-655", "title": "Exact Nash computation in extensive-form games: resolved classification", "jel_code": "C72", "source_id": "OP-0141"}
% !TeX program = xelatex
\documentclass[11pt,a4paper]{article}
\usepackage[margin=23mm]{geometry}
\usepackage{fontspec}
\setmainfont{Latin Modern Roman}
\usepackage{amsmath,amssymb,mathtools}
\usepackage{xcolor,enumitem,booktabs,longtable,array,xurl,hyperref}
\definecolor{muted}{HTML}{53636C}
\hypersetup{colorlinks=true,urlcolor=blue,linkcolor=blue}
\setlength{\parindent}{0pt}
\setlength{\parskip}{5pt}
\newcommand{\R}{\mathbb R}
\newcommand{\E}{\mathbb E}
\newcommand{\N}{\mathbb N}
\newcommand{\ind}{\mathbf 1}
\newcommand{\Var}{\operatorname{Var}}
\newcommand{\Cov}{\operatorname{Cov}}
\newcommand{\supp}{\operatorname{supp}}
\DeclareMathOperator*{\argmin}{arg\,min}
\DeclareMathOperator*{\argmax}{arg\,max}
\newcommand{\entrymeta}[3]{{\sffamily\small\color{muted}\textbf{\texttt{#1}}\quad|\quad #2\par #3\par}}
\newcommand{\Problemref}[1]{\texttt{#1}}
\newcommand{\JELref}[2]{\texttt{#2} (JEL \texttt{#1})}
\begin{document}
\section*{Exact Nash computation in extensive-form games: resolved classification}
\textbf{Problem ID:} \texttt{C72-655}\quad \textbf{JEL:} \texttt{C72}\par
% BEGIN ECONOMICS STATEMENT
\label{problem:OP-0141}
% GAME_THEORY_PROBLEM: {"local_id": "GT-011", "original_number": 11, "kind": "H", "entry_kind": "statement_only", "published": false, "review_required": true, "category": "game_theory"}
\label{game:GT-011}
{\sffamily\small\color{muted}
\textbf{GT-011} | Extensive-form equilibrium complexity\\
\textbf{H} | Historical open classification resolved in the cited paper's 2026 version.
\par}
\subsection*{Definitions and mathematical statement}
An input game is an explicitly listed finite rooted tree. Each nonterminal node is controlled by a player or by chance; chance probabilities and terminal utilities are rational and binary encoded. Each player's decision nodes are partitioned into information sets with identical available actions. Perfect recall means that nodes in one information set have identical sequences of that player's earlier information sets and actions. The number of players is part of the input.
A behavioral strategy $b_i$ specifies a probability distribution at every information set of player $i$. Independent local randomization and chance induce a distribution on terminal nodes and expected payoff $U_i(b)$. The exact Any-Nash search problem is to produce an exact real behavioral profile $b$ such that
\[
 \forall i\ \forall b_i',\qquad U_i(b)\geq U_i(b_i',b_{-i}).
\]
Its general real-search classification is
\[
 \text{Any-Nash in these extensive-form games is }\mathsf{FIXP}\text{-complete}.
\]
This is a reported theorem, rather than the unresolved question asserted in the supplied catalog.
\subsection*{Short English statement}
The general exact extensive-form Nash search problem has a fixed-point complexity classification.
\subsection*{Variants and scope boundaries}
Input size is the explicit tree and information partition, not the exponentially larger strategic-form table. Exact real output is interpreted in the $\mathsf{FIXP}$ search model; arbitrary equilibria need not have rational coordinates. Two-player subclasses, geometric approximations, and imperfect-recall games require separate statements.
\subsection*{Source relationship and status}
The introductory historical discussion in [S1] mentions the old gap. Section 6.2 and Theorem 5 of version 2, dated 29 July 2026, close it for the general class. Reading the historical sentence as the paper's final conclusion caused the original catalog's incorrect open-status claim.
\subsection*{References}
\begingroup\footnotesize\raggedright
{\footnotesize\raggedright
\textbf{[S1]} Vincent Cheval, Florian Horn, Soumyajit Paul, and Mahsa Shirmohammadi (2026). \emph{On the Complexity of the Optimal Correlated Equilibria in Extensive-Form Games}. arXiv:2507.11509v2, 29 July 2026. Locator: Section 2 for the model; Section 6.2 and Theorem 5 for the classification. \url{https://arxiv.org/pdf/2507.11509v2}\par}
\par\endgroup

\subsection*{Common conventions: Source mathematical conventions}

\textbf{Finite games.} For a finite set $A$, $\Delta(A)$ is its probability
simplex. Players choose independent mixed strategies in Nash equilibrium;
correlation is allowed only when explicitly part of the solution concept.
In a game with payoff functions $u_i$ and strategy profile $x$, an additive
$\varepsilon$ Nash equilibrium satisfies
\[
 u_i(x)\geq u_i(x_i',x_{-i})-\varepsilon
 \quad\text{for every player $i$ and every admissible deviation $x_i'$}.
\]
For numerical approximation, payoffs are normalized as locally specified.
An $\varepsilon$ payoff residual is distinct from distance at most
$\varepsilon$ to the set of exact equilibria. Point convergence, convergence
of empirical play, and convergence in distance to an equilibrium set are
also distinct.

\textbf{Computational representations.} Unless stated otherwise, a complexity
question takes rational data with binary encoding; polynomial time is measured
in total bit length. A fixed-accuracy polynomial algorithm is not necessarily
polynomial in $1/\varepsilon$, and neither is necessarily polynomial in
$\log(1/\varepsilon)$. Succinct games and value, demand, or sampling oracles
must declare their interfaces. Exact real solutions require an explicit
representation; an unspecified list of arbitrary real numbers is not a
finite computational output. A claimed algorithmic barrier is meaningful
only for its specified model and reductions.

\textbf{Dynamic games.} Histories, information, observation, and allowable
randomization are part of the model. A stationary strategy depends only on
the current state; a positional strategy in a deterministic graph game
chooses one action at each controlled vertex. Nash equilibrium at an initial
state and equilibrium after every history are different requirements.
An expectation of a pathwise $\liminf$ payoff, a $\liminf$ of expected averages,
a limiting expected average, and a uniform finite-horizon equilibrium are
different criteria. None is silently substituted for another.

\textbf{Allocation and mechanisms.} A complete allocation assigns every item
exactly once unless otherwise stated. Zero-valued goods, negative values,
capacity constraints, feasibility, lotteries, and copies of goods affect
fairness definitions and are specified locally. Dominant-strategy incentive
compatibility, Bayesian incentive compatibility, universal truthfulness,
and truthfulness in expectation impose different inequalities. Whenever
an objective ratio has zero denominator, its convention is stated or those
instances are excluded.

\textbf{Infinite-dimensional models.} Expectations, integrals, stochastic
equations, and optimization objectives require their local regularity,
integrability, filtrations, and admissible domains. A backward--forward
mean-field system is a boundary-value problem; it is not an initial-value
semiflow without an additional construction. Long-time, large-population,
and vanishing-noise limits require an order of limits and a convergence
topology. If a minimum need not be attained, the objective is its infimum
or supremum and the approximation requirement is stated separately.
% END ECONOMICS STATEMENT
\end{document}
